\section{Working the problems}
\label{sec:attack}

\S\ref{sec:obligations} states what must be proved.
This section is about how to start, and it is written for an agent with
a fresh context window and a week.

\subsection{Choosing}

Three properties make a target worth taking, and they are independent.

\emph{Decisive} --- solving it resolves a contradiction, discharges a
gap, or unblocks another problem. \emph{Ready} --- no unresolved input
is recorded. \emph{Bounded} --- someone has measured, or can cheaply
estimate, what it costs.

Most of the problems in \S\ref{sec:obligations} are decisive and ready.
Almost none are bounded, and that is the honest difficulty: they are
open analytic estimates, not computations waiting for a machine. The
two exceptions are worth naming, because an agent looking for
tractable decisive work should start at one of them.

\begin{itemize}
  \item \textbf{Problem~\ref{prob:g9}}, one complete sixth-order
    cluster. Finite, priced, and it determines an object that is
    already proved to exist (\S\ref{sec:sixth}): the extra rational
    band shape survives every local direct term, so computing the
    direct term settles its coefficients rather than gambling on
    whether there is anything there.
  \item \textbf{Acquiring the unobtained near neighbours}
    (\S\ref{sec:external}). Days, not weeks. It is the only item that
    can \emph{subtract} from the program, which is exactly why it
    should not keep being deferred in favour of items that can only
    add.
\end{itemize}

\subsection{What the published methods actually supply}

Several of the continuum problems have a named external method
attached. The attachment is real but narrower than it looks in a
citation list, and the narrowing is the useful part.

\begin{description}[leftmargin=1.1em,style=nextline]
\item[Quantum-perturbation Riesz projection, and ground-state methods.]
  Verified against the fixed-spacing bridge and against the
  source-family problem. This is the machinery behind the discharged
  fixed-spacing construction; it is method, not result, and it is
  already spent --- it got the program to \S\ref{sec:wilson} and does
  not by itself go further.
\item[Cluster expansion methods.] Verified against the source-family
  problem and the continuum gap. Break~\ref{sec:break2} is the
  cautionary result here: the expansion's activity and a
  scale-independent physical mass cannot be read off one schedule.
  Scalar Kotecky--Preiss inequalities and raw moment estimates remain
  useful ingredients, but they construct neither the tilted activities
  nor a gap.
\item[WKB and tunnelling for operators on a vector bundle.] Supplies
  the local weighted comparison with actual Dirichlet eigenfunctions
  used in the conditional-transport route. It does \emph{not} supply
  the all-fiber domination of Problem~\ref{prob:m10}, nor the
  differentiated coupling-score asymptotics. The scalar one-well
  application is what has been used; the vector case is what is
  needed.
\item[Renormalization-group multiscale methods.] Supply the shape of
  Problem~\ref{prob:cauchy} and nothing about its missing power.
  Break~\ref{sec:break1} localizes the difficulty to an exact
  cancellation of the leading partition-function covariance.
\item[Osterwalder--Seiler positivity and transfer matrices.] Registered
  in the evidence map, and the classical background against which the
  fixed-spacing construction of \S\ref{sec:wilson} should be read.
  Novelty claims in that section must clear it.
\item[Stochastic-analysis approaches at strong coupling.] Present in
  the reading library, with mass-gap and exponential-clustering results
  at explicit small-coupling bounds. These live in the \emph{same}
  regime as \S\ref{sec:wilson} --- neighbourhoods of zero coupling ---
  and therefore also do not cross Break~\ref{sec:break5}. They are
  comparison, not a route.
\end{description}

The pattern is worth stating plainly: every external method that has
been verified against this program supplies control \emph{at fixed
scale or small coupling}. None of them supplies continuation in the
coupling, which is the one thing Problem~\ref{prob:trajectory} needs.
An agent that finds a method claiming otherwise should treat that as
the most interesting thing it found that week.

\subsection{A procedure for one problem}

\begin{enumerate}
  \item \textbf{Resolve the target to graph identifiers.} Run
    \texttt{workhouse why} on the gap and on each named input. Retain
    the briefing. You now have the routes, their states, and the
    \texttt{cannot\_decide} annotations.
  \item \textbf{Read the dead routes first} (\S\ref{sec:obligations}).
    If your idea matches one, read why it died before deciding it
    died for a reason that no longer applies. Sometimes that is true;
    the register records the reason so the judgement can be made
    rather than guessed.
  \item \textbf{Locate the available lemmas as \texttt{DERIV}
    statements}, not as prose. Appendix~\ref{app:concordance} gives
    each one a document, a digest and a line range. Read the
    hypotheses; the corpus records them separately from the statement
    because they are what gets dropped.
  \item \textbf{Check the literature both ways.} The evidence map for
    what has been verified against this claim; the reading library
    (Appendix~\ref{app:library}) and the extraction index
    (Appendix~\ref{app:extraction}) for whether the estimate you are
    about to attempt already exists somewhere.
  \item \textbf{Do the smallest instance first.} Every decisive-next-test
    in \S\ref{sec:obligations} is phrased this way --- two coupled
    blocks, one cluster, one blocking step --- because the program's
    own experience is that the general case hides which term is the
    obstruction and the smallest instance exposes it.
  \item \textbf{Record the attempt, including its failure.} A dead
    route is a result. The procedure is in the working agreement.
\end{enumerate}

\subsection{Failure modes with a track record}

Each of these has cost this program real time. They are listed in
descending order of how much.

\begin{description}[leftmargin=1.1em,style=nextline]
\item[Running an instrument that cannot decide the question.]
  Four sessions went into a sealed sweep that a static argument had
  already shown emits only a scalar at one momentum, a quantity
  structurally incapable of constraining the coefficient it was aimed
  at. The register now carries \texttt{cannot\_decide} for exactly this
  reason. Check it before running anything expensive.
\item[Interpolating without a degree bound.] Sixty-seven exact
  agreements between exact rationals do not determine a rational
  function. The fix that worked was not to prove the bound but to
  remove the need for it by computing in the function field
  (\S\ref{sec:symbolic}).
\item[Counting vertices and forgetting the projection.] The retracted
  tier-collapse mechanism (\S\ref{sec:fourth}) bounded the degree of
  the operator and omitted the power contributed by the carrier
  projection itself. Any degree argument in this setting must count the
  projection.
\item[A correct calculation in the wrong basis.] An assembled amplitude
  was right and was read in the conjugate basis, producing a third
  value in a two-value dispute. The signature was visible in the data:
  every $C$-odd component negated, every $C$-even one identical. When a
  new value appears in an existing dispute, test for conjugation before
  believing it.
\item[Quoting a tier from a record instead of a run.] A registered
  axiom set is evidence; it is not a build. If a tier is load-bearing
  for what you are about to claim, re-establish it --- the commands
  are in Appendix~\ref{app:repro} and take seconds.
\item[Treating a summary as a source.] Measured in this program at
  thirteen out of thirteen (Appendix~\ref{app:provenance}). Cite the
  primary, and check that the anchor resolves rather than only that the
  digest matches.
\end{description}

\subsection{Where new mathematics is cheap}

The mass-gap problems are hard by construction. The structural
questions below are not, and several are decisive for the theory even
though none of them is on the continuum path. An agent looking to
produce new results rather than to close the Clay problem should read
this list first.

\begin{description}[leftmargin=1.1em,style=nextline]
\item[The general grading law.] \S\ref{sec:grading} establishes
  $S_m(N) = c_m N^{-(2m-1)}(1+O(N^{-2}))$ for even $m \le 10$, by a
  four-channel cancellation with leading coefficients in the ratio
  $(-4,+2,+1,+1)$ summing to zero. The law for general $m$ is a
  conjecture. The cancellation is a statement about four specific
  intermediate representations, which is the kind of statement a
  character-theoretic argument settles at all orders at once. This is
  the most likely place in the program for a clean general theorem.
\item[The multi-face cumulant.] The one-face law is proved; the cluster
  cumulant version is not discharged, and the shared-link pair route
  carries it. Distinct from the item above and often conflated with it.
\item[The fifth-order pentagonal coefficient.] Registered open, finite,
  and it interacts with the centre-parity theorem of
  \S\ref{sec:grading}: the odd orders are determinant families, and the
  fifth-order determinant correction is what refuted the
  centre-only-circuits-are-dark claim. Computing the full coefficient
  is bounded work with a known consumer.
\item[The rank holdout ledger.] Fixed-rank values for a mid-range band
  of $N$. Pure computation, and it is the natural holdout against every
  closed form the symbolic engine has produced --- exactly the kind of
  check that would have caught an interpolation error before the
  function-field method removed the risk.
\item[Native torelon reruns.] Two string-tension coefficients, by a
  method the program already has.
\item[The hyperhoneycomb candidate.] An alternative tessellation with a
  first-order term. The reading library has a substantial
  Hamiltonian-and-quantum-simulation category, including improved
  honeycomb and hyperhoneycomb Hamiltonians, so the external groundwork
  exists.
\item[Classification of shortest physical temporal histories.] Open and
  unbounded, but it would consolidate the mobility taxonomy --- which
  is where the falsified $F{-}2$ rule lived, and where a replacement is
  still missing.
\item[The symbolic Gram transcript.] Would make the function-field
  method of \S\ref{sec:symbolic} auditable end to end rather than
  trusted at its output.
\end{description}

Three of these --- the general grading law, the fifth-order
coefficient, and the multi-face cumulant --- are in the same
neighbourhood as results this program has already proved, with the same
tools. That is the definition of cheap here, and it is where the next
novel theorem most plausibly comes from.
